\clearpage
\section{Results and Discussion}

\subsection{Benchmark Methodology}

All benchmark values used for this report are available in the repository: \href{https://github.com/pcd-26/PCD-26}{pcd-26/PCD-26}.
The figures in \texttt{assignment-1/report/images/} visualize the same summaries.
The report uses a speedup snapshot with \textbf{1000}, \textbf{1500}, \textbf{2000}, and \textbf{2500} balls at \textbf{600} steps, seed \texttt{42}, and a worker-count snapshot for \textbf{2500} balls at \textbf{1000} steps, seed \texttt{42}.

\begin{itemize}
    \item \textbf{Environment}: OpenJDK 17.0.19+10 on Windows 11, amd64;
    \item \textbf{Visible cores}: 16;
    \item \textbf{Measured metrics}: median elapsed time, median throughput, and speedup;
    \item \textbf{Comparability rule}: only runs with the same workload parameters are compared.
\end{itemize}

The first table summarizes the speedup snapshot; the second table and the second figure summarize the worker-count snapshot.

\begin{itemize}
    \item \textbf{Speedup table}: compares the sequential median with the concurrent medians across workloads;
    \item \textbf{Worker table}: shows the 2500-ball scaling trend as worker count changes;
    \item \textbf{Interpretation}: larger workloads benefit more because the parallelizable physics work outweighs the remaining serial portion.
\end{itemize}

The speedup snapshot contains the following values:

\begin{table}[htbp]
    \centering
    \fontsize{12}{16}\selectfont
    \setlength{\tabcolsep}{4pt}
    \renewcommand{\arraystretch}{0.9}
    \begin{tabular}{rlrrr}
        \hline
        Balls & Implementation & Seq. median [ms] & Concurrent median [ms] & Speedup \\
        \hline
        1000 & threads  & 129.4819 & 117.7674 & 1.099 \\
        1000 & executor & 129.4819 & 107.5256 & 1.204 \\
        1500 & threads  & 207.5886 & 145.7583 & 1.424 \\
        1500 & executor & 207.5886 & 148.2290 & 1.400 \\
        2000 & threads  & 323.5650 & 182.7458 & 1.771 \\
        2000 & executor & 323.5650 & 189.7117 & 1.706 \\
        2500 & threads  & 433.9208 & 230.3475 & 1.884 \\
        2500 & executor & 433.9208 & 232.9631 & 1.863 \\
        \hline
    \end{tabular}
    \caption{Median speedup values for the local speedup snapshot (1000, 1500, 2000, and 2500 balls; 600 steps; seed 42) on Windows 11 amd64 with OpenJDK 17.0.19+10 and 16 visible cores.}
\end{table}

\begin{figure}[htbp]
    \centering
    \includegraphics[width=\linewidth]{images/speedup-vs-balls}
    \caption{Speedup as the number of balls increases for the local speedup snapshot (1000 to 2500 balls; 600 steps; seed 42) with 16 workers.}
\end{figure}

The worker-count snapshot contains the medians for the 2500-ball scalability workload with \textbf{1000} steps and seed \texttt{42}.

\begin{table}[htbp]
    \centering
    \fontsize{12}{16}\selectfont
    \setlength{\tabcolsep}{4pt}
    \renewcommand{\arraystretch}{0.9}
    \begin{tabular}{lrrr}
        \hline
        Implementation & Workers & Median [ms] & Speedup vs. seq. \\
        \hline
        threads  & 1  & 1191.1281 & 0.612 \\
        threads  & 2  & 727.4179  & 1.003 \\
        threads  & 4  & 462.9961  & 1.575 \\
        threads  & 8  & 367.9059  & 1.982 \\
        threads  & 16 & 383.0851  & 1.904 \\
        threads  & 17 & 394.2740  & 1.850 \\
        executor & 1  & 1179.4044 & 0.618 \\
        executor & 2  & 726.8271  & 1.003 \\
        executor & 4  & 496.5341  & 1.469 \\
        executor & 8  & 384.8939  & 1.895 \\
        executor & 16 & 387.3602  & 1.883 \\
        executor & 17 & 415.1978  & 1.756 \\
        \hline
    \end{tabular}
    \caption{
        Worker-count medians for the local 2500-ball scalability workload
        (1000 steps; seed 42) comparing the concurrent runs against the
        sequential baseline at 1, 2, 4, 8, 16, and 17 workers.
    }
\end{table}

The \textbf{benchmark evidence} shows a real but \textbf{sublinear} gain.
On the speedup snapshot, both concurrent versions beat the sequential baseline at every measured workload.
The executor-based engine is ahead at 1000 balls, while the threaded engine is ahead from 1500 through 2500 balls and reaches the largest speedup, \texttt{1.883766}.
On the worker-count snapshot, both versions improve through 8 workers, where they reach their best medians, then regress at 16 and 17 workers.
This saturation is consistent with the remaining \textbf{serial bottleneck}, scheduling costs, and contention once the useful parallelism of this workload is exhausted.

\begin{figure}[htbp]
    \centering
    \includegraphics[width=0.84\linewidth]{images/speedup-vs-workers}
    \caption{Speedup as the worker count changes for the local 2500-ball scalability workload (1000 steps; seed 42).}
\end{figure}
